\uuid{q5lg}
\exo7id{3996}
\titre{exo7 3996}
\auteur{quercia}
\organisation{exo7}
\datecreate{2010-03-11}
\isIndication{false}
\isCorrection{false}
\chapitre{Dérivabilité des fonctions réelles}
\sousChapitre{Fonctions convexes}
\cpge{mpsi-mp2i.convexite.b, mpsi-mp2i.derivabilite.c, pcsi.derivabilite.c, ptsi.derivabilite.c, mpsi-mp2i.fonctions-valeurs-reelles-complexes.b, pcsi.techniques-calcul-fonctions.b, ptsi.techniques-calcul-fonctions.b, bcpst1.derivation.b, bcpst1.suites-reelles}
\module{Analyse}
\niveau{L1}
\difficulte{}

\contenu{
\texte{
Soit $f : {[0,+\infty[} \to {\R}$ concave, dérivable, croissante.
}
\begin{enumerate}
    \item \question{Montrer que : $\forall\ x\ge 1,\ f(x+1)-f(x) \le f'(x) \le f(x)-f(x-1)$.}
    \item \question{On pose : $\begin{cases} u_n = f'(1) + f'(2) + \dots + f'(n) - f(n) \cr
                       v_n = f'(1) + f'(2) + \dots + f'(n) - f(n+1).\cr\end{cases}$
    Montrer que ces suites convergent.}
    \item \question{On prend $f(x) = \ln x$. Soit $\gamma = \lim_{n\to\infty} u_n$
    (constante d'Euler). Calculer $\gamma$ à $10^{-2}$ près.}
\end{enumerate}
}
